\documentclass[11pt,a4paper]{article}

\usepackage{fontspec}
\usepackage{polyglossia}
\setdefaultlanguage{russian}
\setotherlanguage{english}
\setmainfont{Arial}
\setsansfont{Arial}
\setmonofont{Consolas}
\usepackage{unicode-math}
\setmathfont{Latin Modern Math}

\usepackage[a4paper,top=18mm,bottom=20mm,left=17mm,right=17mm]{geometry}
\usepackage{microtype}
\usepackage{amsmath}
\usepackage{booktabs}
\usepackage{array}
\usepackage{tabularx}
\usepackage{longtable}
\usepackage{pdflscape}
\usepackage{enumitem}
\usepackage{xcolor}
\usepackage{hyperref}
\usepackage{fancyhdr}
\usepackage{titlesec}
\usepackage{lastpage}

\definecolor{navy}{HTML}{153451}
\definecolor{blue}{HTML}{2A6F97}
\definecolor{cyan}{HTML}{4A9BB8}
\definecolor{lightblue}{HTML}{EAF3F7}
\definecolor{lightgreen}{HTML}{EDF5EE}
\definecolor{lightred}{HTML}{F8ECEA}
\definecolor{lightamber}{HTML}{FAF3E1}
\definecolor{textgray}{HTML}{4D5B66}
\definecolor{linegray}{HTML}{D8E0E5}

\hypersetup{
  colorlinks=true,
  linkcolor=navy,
  urlcolor=blue,
  pdftitle={WiRE-TPP: протокол экспериментов, corrected all-12 и алгоритм обучения},
  pdfauthor={Экспериментальный протокол WiRE-TPP},
  pdfsubject={Signed-Wishart random effects for temporal point-process clustering},
  pdfkeywords={WiRE-TPP, Signed-Wishart, THP, COTIC, clustering, all-12}
}

\pagestyle{fancy}
\fancyhf{}
\fancyhead[L]{\small\color{textgray}WiRE-TPP / Signed-Wishart}
\fancyhead[R]{\small\color{textgray}Экспериментальный протокол / 07.08.2026}
\fancyfoot[L]{\small\color{textgray}Corrected all-12, абляции и отрицательные результаты}
\fancyfoot[R]{\small\color{textgray}\thepage\ / \pageref{LastPage}}
\renewcommand{\headrulewidth}{0.4pt}
\renewcommand{\footrulewidth}{0pt}
\setlength{\headheight}{14pt}

\titleformat{\section}{\Large\bfseries\color{navy}}{\thesection}{0.65em}{}
\titleformat{\subsection}{\large\bfseries\color{blue}}{\thesubsection}{0.65em}{}
\titleformat{\subsubsection}{\normalsize\bfseries\color{cyan}}{\thesubsubsection}{0.65em}{}
\titlespacing*{\section}{0pt}{1.25em}{0.5em}
\titlespacing*{\subsection}{0pt}{0.95em}{0.35em}
\titlespacing*{\subsubsection}{0pt}{0.75em}{0.25em}

\setlist[itemize]{leftmargin=1.45em,itemsep=0.18em,topsep=0.3em}
\setlist[enumerate]{leftmargin=1.65em,itemsep=0.22em,topsep=0.3em}
\setlength{\parindent}{0pt}
\setlength{\parskip}{0.44em}
\setlength{\emergencystretch}{3em}
\renewcommand{\arraystretch}{1.14}

\newcommand{\code}[1]{\texttt{#1}}
\newcommand{\R}{\mathbb{R}}
\newcommand{\Ktarget}{K^{\star}}
\newcommand{\boxnote}[2]{%
  \begin{center}\setlength{\fboxsep}{9pt}%
  \colorbox{#1}{\begin{minipage}{0.94\textwidth}\raggedright #2\par\end{minipage}}
  \end{center}}
\newcommand{\verdict}[1]{\boxnote{lightred}{\textbf{Итог.} #1}}
\newcommand{\rationale}[1]{\boxnote{lightgreen}{\textbf{Обоснование.} #1}}
\newcommand{\caveat}[1]{\boxnote{lightamber}{\textbf{Ограничение.} #1}}
\newcommand{\sourcepath}[1]{\par\noindent\textcolor{textgray}{\footnotesize\path{#1}}\par}

\begin{document}

\begin{titlepage}
  \thispagestyle{empty}
  \vspace*{11mm}
  {\sffamily\bfseries\fontsize{27}{32}\selectfont\color{navy}
    WiRE-TPP / Signed-Wishart\par}
  \vspace{3mm}
  {\sffamily\fontsize{19}{24}\selectfont\color{blue}
    Протокол экспериментальной программы,\par
    corrected all-12 и полный алгоритм обучения\par}
  \vspace{8mm}
  {\color{cyan}\rule{\textwidth}{1.6pt}}
  \vspace{8mm}

  {\large Сводка воспроизводимости, последовательных абляций THP/COTIC,
  matched-DGP экспериментов и отрицательных результатов с точной фиксацией
  фактически применённых метапараметров.}

  \vspace{7mm}
  \verdict{Совокупность экспериментов не подтверждает устойчивого общего
  преимущества Signed-Wishart по кластеризации. На corrected all-12 COTIC
  остаётся сильным и слегка выигрывает у no-W в среднем, тогда как THP в среднем
  немного проигрывает. На специально согласованном Wishart-DGP обнаружен слабый
  средний эффект, но его знак меняется между seeds. Обучаемая $\nu$ и увеличение
  числа train-MC samples проблему не устранили.}

  \vspace{4mm}
  \boxnote{lightblue}{
    \textbf{Правило чтения результатов.} Выбор checkpoint и конфигурации
    производится по максимальной validation purity, затем по меньшему validation
    NLL. Test используется только описательно после выбора. Строки, где решение
    было принято вручную или после просмотра test, прямо помечены и не считаются
    подтверждающим сравнением.
  }

  \vfill
  \begin{tabular}{@{}ll@{}}
    \textbf{Дата фиксации:} & 7 августа 2026 г. \\
    \textbf{Рабочая директория:} & \code{C:\textbackslash Users\textbackslash missi\textbackslash Documents\textbackslash WiRETPP} \\
    \textbf{GPU:} & NVIDIA GeForce RTX 4060 Laptop, 8188 MiB \\
    \textbf{Основные backbones:} & THP и COTIC \\
    \textbf{Главная задача:} & кластеризация marked event sequences \\
  \end{tabular}
  \vspace{8mm}
\end{titlepage}

\begingroup
\footnotesize
\setlength{\parskip}{0pt}
\tableofcontents
\endgroup
\newpage

\section{Краткое резюме и статус доказательств}

\subsection{Что действительно установлено}

\begin{enumerate}
  \item Репозиторий и опубликованные checkpoints воспроизводимы на уровне
  inference: повторный MC-аудит фиксированных весов дал отклонения suffix NLL
  не более $6.1\cdot10^{-4}$ и purity не более $1.88\cdot10^{-3}$.
  \item В ранней 300-epoch реконструкции Wishart улучшал predictive suffix NLL
  во всех шести paired-сравнениях, но clustering metrics заметно менялись при
  retraining. Значит, MC evaluation был стабилен, а optimization path --- нет.
  \item На corrected all-12, объединённом по seeds 0, 1, 2, COTIC-Wishart имеет
  среднюю test purity $0.98874$ против $0.98715$ у COTIC no-W
  ($+0.00159$). THP-Wishart имеет $0.97245$ против $0.97677$
  ($-0.00432$).
  \item Самый трудный случай --- \code{sin\_K5\_C5} с THP. На трёх seeds
  THP-Wishart даёт $0.8667\pm0.0900$ test purity, no-W ---
  $0.9250\pm0.0260$. COTIC на том же датасете остаётся почти идеальным:
  $0.9883\pm0.0076$ для Wishart.
  \item Заморозка encoder во время overcomplete elimination и отдельная
  K=1-инициализация улучшали отдельные односидовые точки, но эффект сильно
  зависел от pretraining length и не переносился как единый устойчивый рецепт.
  \item Matched block-local Wishart-DGP с gain 3 дал среднюю validation purity
  $0.8417$ для Wishart против $0.8250$ для no-W, но paired deltas равны
  $+0.125,-0.025,-0.050$. Средний знак положителен, устойчивого эффекта нет.
\end{enumerate}

\subsection{Что не следует утверждать}

\begin{itemize}
  \item Нельзя утверждать, что latent Wishart всегда улучшает кластеризацию:
  aggregate THP и несколько сильных DGP-сценариев это опровергают.
  \item Нельзя считать \code{lal\_bos} полным воспроизведением опубликованного
  LaL. Это только fixed-$K$ BOS-control с общим COTIC; random-walk, split/merge/
  delete и EM-процедура LaL в нём отсутствуют.
  \item Нельзя выбирать $\eta_{\alpha}=10^{-4}$ по test. В beta$_0=1$ sweep
  validation winner был $10^{-5}$; $10^{-4}$ был затем зафиксирован как
  инженерное решение для следующей постановки.
  \item Нельзя интерпретировать движение $\alpha$ или $\nu$ как доказательство
  идентифицированного random effect без paired improvement в held-out metrics.
\end{itemize}

\section{Карта экспериментальной программы}

\begin{longtable}{p{0.12\textwidth}p{0.24\textwidth}p{0.16\textwidth}p{0.36\textwidth}}
\toprule
Дата & Эксперимент & Статус & Главный результат \\
\midrule
\endfirsthead
\toprule
Дата & Эксперимент & Статус & Главный результат \\
\midrule
\endhead
01.08 & Published-checkpoint audit и полный 300-epoch retraining & завершён & Fixed-weight inference воспроизводим; Wishart улучшает NLL, clustering retraining нестабилен. \\
01.08 & Signed K3C5 DGP, 3 seeds & завершён & Moderate heterogeneity: signed лучше $\alpha=0$; strong: signed хуже на всех seeds. \\
02.08 & THP \code{sin\_K5\_C5}, fixed split, 5 seeds & завершён & Wishart $+0.019$ test purity в среднем, но сохраняет four-mode collapse. \\
02.08 & Optuna, frozen exploration, MC и scheduler & завершён & Нет устойчивой точки; больше MC не помогает; beta curriculum даёт лишь $+0.003$ validation purity. \\
03.08 & Two-phase COTIC \code{K5\_C5} & завершён, 1 seed & K=1 pretrain + freeze + DPP дал val/test $0.995/1.000$; требуется multiseed. \\
04--05.08 & Original и corrected all-12 & завершён & 12 datasets $\times$ 2 architectures $\times$ 2 variants $\times$ 3 seeds; COTIC стабилен, THP \code{sin\_K5} нестабилен. \\
06.08 & Frozen backbone LR $\times$ warmup, 25 jobs & завершён & Лучший seed-0 val/test $0.945/0.910$ при $\eta_{bb}=3\cdot10^{-4}$, warmup 50. \\
06.08 & $\eta_{\alpha}\times$ K=1 pretrain, 25 jobs & завершён & Лучший val $0.935$ при pretrain 300 и $\eta_{\alpha}=10^{-5}$; test $0.890$ описательно. \\
07.08 & Big-protocol beta$_0=1$, $\eta_{\alpha}$ sweep, 5 jobs & завершён & Validation winner $10^{-5}$: val/test $0.910/0.885$; $10^{-4}$: $0.905/0.925$. \\
07.08 & Overcomplete backbone LR sweep, 7 planned & частичный, 5/7 & Из завершённых лучший $10^{-6}$: val/test $0.870/0.835$; прогон остановлен. \\
07.08 & Первый full matched Wishart-RE suite, 9 planned & остановлен, 2/9 & Обе модели collapsed в один cluster; постановка признана дефектной. \\
07.08 & Исправленный block-local DGP gain 1/3/5 & debug завершён & Gain 3 дал слабый средний Wishart эффект, но paired sign нестабилен. \\
07.08 & Fixed/learned $\nu$, LR$_\nu$, MC 1/2/4 & debug завершён & Fixed $\nu=25$ выбран validation rule; learnable $\nu$ и больше MC purity не улучшили. \\
\bottomrule
\end{longtable}

\caveat{В workspace есть дополнительные smoke, unit и аварийные каталоги.
Они используются для аудита реализации, но не включаются в количественный
benchmark. Статусы \code{running} в некоторых \code{suite\_progress.json}
являются stale metadata после ручной остановки и не означают живой процесс.}

\section{Данные, split и метрики}

Траектория записывается как
\[
x_i=\{(t_{ij},c_{ij})\}_{j=1}^{n_i},\qquad c_{ij}\in\{1,\ldots,C\},
\qquad t_{ij}\le T_i.
\]
Во всех DAN all-12 экспериментах $C=5$. Dataset name кодирует истинное число
кластеров: \code{K2}, \code{K3}, \code{K4}, \code{K5}. Рассматриваются три
семейства: обычные Hawkes-процессы \code{K*\_C5}, синусоидальные
\code{sin\_K*\_C5} и усечённые \code{trunc\_K*\_C5}.

Для all-12 и поздних sweeps используется один и тот же split:
\code{numpy.RandomState(42).shuffle}, затем 80/10/10. Model seeds меняют
инициализацию, minibatches и Wishart sampling, но не split. Для
\code{sin\_K5\_C5} это 1600/200/200 траекторий.

\begin{center}
\begin{tabularx}{\textwidth}{>{\bfseries}p{0.18\textwidth}X X}
\toprule
Множество & Разрешено & Запрещено \\
\midrule
Train & gradients, OT dual, hyperprior, pruning score & использование labels в objective \\
Validation & purity/NLL checkpoint selection & parameter fitting и test-dependent rule \\
Test & один descriptive pass после выбора & выбор checkpoint, LR, $\alpha$, $\nu$ или DGP \\
\bottomrule
\end{tabularx}
\end{center}

Основная метрика model selection --- validation purity; при равенстве purity
выбирается меньший validation NLL per exposure. Test-метрики: purity, ARI, NMI,
NLL per exposure. Purity допускает many-to-one matching и поэтому дополняется
ARI/NMI и размерами кластеров.

\rationale{Purity была выбрана для сопоставимости с LaL. ARI и NMI нужны,
поскольку высокая purity может скрывать дублирование одной истинной моды двумя
предсказанными компонентами. NLL проверяет predictive fit, но не заменяет
clustering metric: во многих runs NLL улучшался одновременно с ухудшением
purity.}

\section{Модель и полный алгоритм обучения}

\subsection{Базовый mixture TPP}

Shared encoder $H_{\theta}$ строит историю, а component head $G_{\phi_k}$
задаёт $C$ положительных intensities $\lambda_{kc}^{(0)}(t\mid H_{\theta})$.
Для компоненты $k$:
\[
\ell_{ik}=\sum_j\log\lambda_{k,c_{ij}}(t_{ij})-
\int_0^{T_i}\sum_{c=1}^{C}\lambda_{kc}(t)\,dt.
\]
Backbone может быть THP или COTIC; Wishart layer не меняет его контракт.

\subsection{K=1 pretraining и expansion}

Сначала одна компонента обучается по обычному TPP NLL. Затем её head копируется
в $K_0=\Ktarget+5$ блоков:
\[
\phi_k=\phi_1+\varepsilon_k,\qquad
\varepsilon_k\sim\mathcal N(0,\sigma_{exp}^2I).
\]
В all-12 pretraining длится 10 optimizer steps, $\sigma_{exp}=0.01$.

\rationale{K=1 отделяет обучение общей временной динамики от allocation
компонент. Малый шум разрушает permutation symmetry, не уничтожая общий
масштаб intensity. Overcomplete запас из пяти heads даёт backward elimination
несколько кандидатов на каждую финальную компоненту.}

\subsection{Latent Wishart random effect}

При текущем $K$ размер latent matrix равен $D=KC$. Обучаемая средняя матрица
задаётся Cholesky-параметром:
\[
\Omega=LL^{\mathsf T},\qquad \operatorname{tr}\Omega=D,
\qquad \operatorname{diag}(L)=\operatorname{softplus}(r)+10^{-5}.
\]
Траекторный random effect:
\[
W_i\sim\operatorname{Wishart}_{D}(\nu,\Omega/\nu),
\qquad \mathbb E[W_i]=\Omega.
\]
В all-12 $\nu(K)=\lfloor1.5KC\rfloor$: для \code{sin\_K5} это 75 при
$K=10$ и 37 при $K=5$.

\rationale{Одна и та же $W_i$ используется для route и intensity transform на
всей траектории. Trace normalization удаляет неидентифицируемый общий scale,
Cholesky гарантирует SPD, а изменение $\nu$ вместе с $D$ поддерживает примерно
сопоставимый уровень sampling noise при pruning.}

\subsection{Frozen exploration Wishart}

Во время train независимо сэмплируется
$W_{exp}\sim\operatorname{Wishart}(\nu_e,I/\nu_e)$ и применяется curriculum
\[
W_t=(1-\beta_t)W_{learn}+\beta_tW_{exp},
\]
\[
x_t=\frac{t-1}{T-1},\qquad
\beta_t=\beta_0\frac{e^{-d_\beta x_t}-e^{-d_\beta}}
{1-e^{-d_\beta}}.
\]
В all-12 $\beta_0=0.9$, $d_\beta=5$, $\nu_e=200$; в позднем sweep
$\beta_0=1$. На validation, pruning и test всегда $\beta=0$.

\rationale{Curriculum должен задержать раннюю фиксацию случайного basin и затем
точно вернуть inference к learned law. Эксперименты показывают, что сам по себе
этот механизм даёт лишь малый и нестабильный выигрыш.}

\subsection{Signed convex intensity transform}

Из $W$ строится correlation matrix и off-diagonal coupling:
\[
R=\operatorname{diag}(W)^{-1/2}W\operatorname{diag}(W)^{-1/2},
\qquad B=R-I.
\]
Пусть $\lambda^{(0)}\in\R_+^{KC}$ --- base intensities,
$a=\lambda^{(0)}/(\epsilon+\sum_d\lambda_d^{(0)})$ и
$q=\operatorname{softplus}^{-1}(\lambda^{(0)})$. Тогда
\[
p=Ba,\qquad
\widetilde q=(1-\alpha)q+\alpha p,
\qquad
\widetilde\lambda=\operatorname{softplus}(\widetilde q),
\qquad \alpha\in[0,1].
\]
В all-12 $\alpha_0=0.1$, projected parameterization, $\eta_\alpha=10^{-4}$.

\rationale{Correlation normalization отделяет interaction structure от scale;
$B=R-I$ удаляет self-effect. Exact inverse softplus сохраняет исходную TPP
интенсивность при $\alpha=0$. Projected $[0,1]$ ограничивает деформацию, а
финальный softplus гарантирует положительность.}

\subsection{Wishart gate и MC-маргинализация}

Для sample $s$ cluster gate равен нормированной диагональной массе его
$C\times C$ блока:
\[
g_{isk}=\log\frac{\sum_{c=1}^{C}(W_{is})_{kc,kc}}
{\sum_{r=1}^{K}\sum_{c=1}^{C}(W_{is})_{rc,rc}}.
\]
Wishart интегрируется, сохраняя component index:
\[
S_{ik}=\log\sum_{s=1}^{S}\exp(\ell_{isk}+g_{isk})-\log S.
\]

\rationale{Gate и intensity используют одну latent $W_i$, поэтому route и
random effect согласованы. Log-sum-exp даёт стабильную дифференцируемую
аппроксимацию marginal likelihood.}

\subsection{Global unbalanced OT dual}

При target marginal $q_k=1/K$, temperature $\tau$, penalty $\rho$ и dual
$u\in\R^K$:
\[
r_{ik}=\operatorname{softmax}_k\left(\frac{S_{ik}-u_k}{\tau}\right),
\qquad
\bar q_k(u)=\operatorname{softmax}_k\left(\log q_k+\frac{u_k}{\rho}\right),
\]
\[
\mathcal F=-\tau\frac1B\sum_i\log\sum_k e^{(S_{ik}-u_k)/\tau}
-\rho\log\sum_k e^{\log q_k+u_k/\rho}.
\]
Primal minimизирует $\mathcal F$, dual максимизирует; после update
$u\leftarrow u-\operatorname{mean}(u)$. В all-12 $\rho=5$ постоянно,
$\eta_u=0.01$ для Wishart и $0.05$ для исторического no-W control.

\rationale{Global dual мягко препятствует component collapse, не навязывая
равный баланс внутри каждого microbatch. Центрирование устраняет additive
gauge. Разный dual LR у W/no-W является историческим confound и учитывается в
интерпретации paired results.}

\subsection{Hyperprior, DPP и полный loss}

Mean-matrix penalty:
\[
\mathcal L_{\Omega}=\frac{\gamma_\Omega}{2N_{tr}}
(\operatorname{tr}\Omega-\log\det\Omega-D).
\]
Functional DPP, если включён:
\[
K^{DPP}_{kr}=\exp\left(-\frac{\|v_k-v_r\|^2}{2h^2}\right)
+\epsilon_{DPP}\delta_{kr},\qquad
\mathcal L_{DPP}=-\frac1K\log\det K^{DPP}.
\]
Итоговая primal objective:
\[
\mathcal L=\mathcal F+\mathcal L_{\Omega}+w_{DPP}(t)\mathcal L_{DPP}.
\]
В all-12 $\gamma_\Omega=1.566$, а $w_{DPP}=0$. DPP применялся только в
поздних frozen-encoder постановках.

\rationale{Hyperprior одновременно является SPD barrier и shrinkage к identity.
DPP нужен после cloning для функционального разведения heads, но после выбора
target $K$ repulsion может мешать likelihood fine-tuning; поэтому в двухфазном
протоколе он затухает до нуля.}

\subsection{Physical backward elimination}

После warmup 300 выполняются пять удалений на шагах 300, 400, 500, 600, 700.
На всех train trajectories с $S_{prune}$ и $\beta=0$ считаются full score
$L_{full}$ и score физически уменьшенной копии $L_{-j}$:
\[
\Delta_j=L_{-j}-L_{full},\qquad j^\star=\arg\min_j\Delta_j.
\]
Удаляются head $j^\star$ и его $C$ координат в $\Omega$; сохраняется principal
submatrix, затем trace нормируется до нового $D$. Primal optimizer
пересоздаётся, dual обнуляется и меняет размерность, lineage обновляется.

\rationale{Full-train marginal NLL измеряет непосредственный ущерб без labels и
validation leakage. Common random numbers уменьшают variance разности.
Physical removal, в отличие от masking, уменьшает и head, и latent dimension;
optimizer moments после смены геометрии больше невалидны.}

\subsection{Один optimizer step}

\begin{enumerate}
  \item Вычислить $\beta_t$, $\rho_t$, DPP weight и, если применимо, encoder LR.
  \item Взять $A$ microbatches по $B$ trajectories; effective batch $BA$.
  \item Для каждой траектории сэмплировать $S_{train}$ matrices, построить signed
  intensities, gates и $S_{ik}$.
  \item Посчитать $\mathcal F$, hyperprior и DPP; вызвать backward от
  $\mathcal L/A$.
  \item Clip global primal norm до 20, сделать Adam step, project $\alpha$.
  \item Сделать maximizing Adam step dual и центрировать $u$.
  \item На pruning step выполнить physical removal и reset optimizers.
  \item Оценить validation с $\beta=0$; checkpoint eligible только при
  $K=\Ktarget$.
\end{enumerate}

\subsection{Validation selection и test}

Лучший eligible checkpoint максимизирует validation purity; tie-break ---
меньший validation NLL. После train загружаются model state и dual лучшего
checkpoint. Test оценивается один раз с $\beta=0$ и $S_{test}=16$ без updates.

\section{Corrected all-12: точный протокол}

\subsection{Состав benchmark}

\begin{itemize}
  \item 12 datasets: \code{K2--K5}, \code{sin\_K2--sin\_K5},
  \code{trunc\_K2--trunc\_K5}; везде $C=5$.
  \item 2 architectures: THP и COTIC.
  \item 2 variants: no-W и Signed-Wishart.
  \item 3 model seeds: 0, 1, 2; split seed всегда 42.
  \item Всего в объединённом анализе 144 runs. Corrected seed-0 Wishart suite
  содержит 24/24; seeds 1--2 suite --- 96/96; seed-0 no-W взят из предыдущего
  завершённого suite.
\end{itemize}

\subsection{Гиперпараметры Wishart all-12}

\begin{longtable}{p{0.30\textwidth}p{0.21\textwidth}p{0.39\textwidth}}
\toprule
Параметр & Значение & Комментарий \\
\midrule
\endfirsthead
\toprule
Параметр & Значение & Комментарий \\
\midrule
\endhead
$K_0,\Ktarget$ & $\Ktarget+5,\Ktarget$ & Всегда ровно пять eliminations. \\
K=1 pretraining & 10 steps & Seed 0 W построил собственный K=1; seeds 1--2 W переиспользовал paired no-W checkpoint byte-identically. \\
Warmup / prune interval / post-prune & 300 / 100 / 600 & Prune steps 300, 400, 500, 600, 700; total 1300. \\
Head clone noise & 0.01 & Symmetry breaking. \\
Neural LR & $10^{-3}$ & Фактический общий optimizer для shared backbone и heads. \\
Historical backbone metadata LR & $10^{-5}$ & Записан в config/result, но не создаёт отдельную group; history логирует $10^{-3}$. \\
Omega LR / alpha LR & $10^{-2}$ / $10^{-4}$ & Без weight decay. \\
Dual LR / $\rho$ & $10^{-2}$ / 5 & $\rho$ постоянно 5. \\
Weight decay / global clip & $10^{-5}$ / 20 & Weight decay только neural parameters. \\
Physical batch / accumulation & 16 / 8 & Effective batch 128. \\
Eval batch & 16 & Validation на каждом optimizer step. \\
MC train/val/test/prune & 2 / 8 / 16 / 8 & Independent fixed-seed RNG streams. \\
$\alpha_0$, parameterization, mode & 0.1, projected, convex & $\alpha\in[0,1]$, temperature 1. \\
$\nu(K)$ & $\lfloor1.5KC\rfloor$ & Integer fixed degrees of freedom. \\
$\beta_0$, decay, $\nu_e$ & 0.9, 5, 200 & Exponential-to-zero; eval всегда $\beta=0$. \\
Mean hyperprior & 1.566 & Trace/logdet penalty. \\
DPP & 0 & Bandwidth/jitter metadata неактивны. \\
Encoder freeze / ramp & false / неактивен & Это не frozen protocol 06.08. \\
Selection & max val purity, min val NLL & Test не участвует. \\
\bottomrule
\end{longtable}

\caveat{No-W control не является идеально однофакторной абляцией Wishart:
его dual LR равен 0.05, physical batch/accumulation 64/2, а seed-0 correction
использует другой K=1 checkpoint. Effective batch в обоих variants равен 128,
а seeds 1--2 имеют exact paired K=1 reuse.}

\subsection{Seed 0: полный corrected Wishart результат}

В следующих таблицах W --- corrected Wishart, no-W --- исторический paired
control. Значения test описательны; LaL --- опубликованный mean из локального
reference-файла:
\sourcepath{artifacts/dan_all_arch_newwishart_beta_cotic7_seed0_20260802/PUBLISHED_DAN_PURITY.csv}.

\begin{landscape}
\scriptsize
\begin{longtable}{llrrrrrrrr}
\toprule
Dataset & Arch & Val W & Test W & ARI W & NLL W & Test no-W & $\alpha$ & Step & LaL \\
\midrule
\endfirsthead
\toprule
Dataset & Arch & Val W & Test W & ARI W & NLL W & Test no-W & $\alpha$ & Step & LaL \\
\midrule
\endhead
K2\_C5 & COTIC & 1.000 & 1.000 & 1.000 & 1.9915 & 1.000 & 0.1079 & 747 & 0.970 \\
K3\_C5 & COTIC & 1.000 & 1.000 & 1.000 & 1.9742 & 1.000 & 0.1077 & 800 & 0.850 \\
K4\_C5 & COTIC & 1.000 & 1.000 & 1.000 & 1.8997 & 1.000 & 0.1165 & 1095 & 0.900 \\
K5\_C5 & COTIC & 1.000 & 1.000 & 1.000 & 1.9231 & 0.965 & 0.1083 & 1283 & 0.840 \\
sin\_K2\_C5 & COTIC & 1.000 & 1.000 & 1.000 & 3.4839 & 1.000 & 0.1097 & 721 & 0.990 \\
sin\_K3\_C5 & COTIC & 1.000 & 0.9917 & 0.9759 & 0.6453 & 0.9917 & 0.1034 & 1258 & 0.990 \\
sin\_K4\_C5 & COTIC & 0.9813 & 0.9438 & 0.8484 & 0.8200 & 0.9500 & 0.1030 & 1001 & 0.920 \\
sin\_K5\_C5 & COTIC & 0.990 & 0.980 & 0.9474 & 3.0326 & 0.990 & 0.1117 & 1281 & 0.920 \\
trunc\_K2\_C5 & COTIC & 1.000 & 1.000 & 1.000 & 3.7149 & 1.000 & 0.1336 & 702 & 1.000 \\
trunc\_K3\_C5 & COTIC & 0.975 & 0.9667 & 0.9040 & 4.2103 & 0.9833 & 0.1455 & 825 & 0.960 \\
trunc\_K4\_C5 & COTIC & 0.9938 & 0.9875 & 0.9655 & 3.8171 & 0.9875 & 0.1383 & 706 & 0.990 \\
trunc\_K5\_C5 & COTIC & 0.950 & 0.975 & 0.9379 & 3.9446 & 0.970 & 0.1395 & 713 & 0.940 \\
\midrule
K2\_C5 & THP & 1.000 & 1.000 & 1.000 & 2.4279 & 1.000 & 0.1753 & 1288 & 0.970 \\
K3\_C5 & THP & 1.000 & 1.000 & 1.000 & 2.3733 & 1.000 & 0.1683 & 1288 & 0.850 \\
K4\_C5 & THP & 1.000 & 0.9938 & 0.9844 & 2.2501 & 0.9875 & 0.1814 & 1274 & 0.900 \\
K5\_C5 & THP & 0.995 & 0.980 & 0.9521 & 2.4411 & 0.990 & 0.1778 & 1223 & 0.840 \\
sin\_K2\_C5 & THP & 0.975 & 1.000 & 1.000 & 3.6827 & 1.000 & 0.1099 & 1300 & 0.990 \\
sin\_K3\_C5 & THP & 1.000 & 0.9833 & 0.9524 & 1.0541 & 0.9917 & 0.1055 & 1280 & 0.990 \\
sin\_K4\_C5 & THP & 0.9063 & 0.9063 & 0.7774 & 1.1977 & 0.900 & 0.1026 & 1265 & 0.920 \\
sin\_K5\_C5 & THP & 0.865 & 0.870 & 0.7166 & 3.3949 & 0.910 & 0.1390 & 1267 & 0.920 \\
trunc\_K2\_C5 & THP & 1.000 & 1.000 & 1.000 & 3.6242 & 1.000 & 0.1528 & 797 & 1.000 \\
trunc\_K3\_C5 & THP & 0.975 & 0.975 & 0.9269 & 4.1716 & 0.9667 & 0.1588 & 858 & 0.960 \\
trunc\_K4\_C5 & THP & 0.9938 & 0.9875 & 0.9655 & 3.8055 & 0.9938 & 0.1719 & 1001 & 0.990 \\
trunc\_K5\_C5 & THP & 0.950 & 0.975 & 0.9399 & 3.9370 & 0.980 & 0.1626 & 851 & 0.940 \\
\bottomrule
\end{longtable}
\end{landscape}

Для seed-0 THP \code{sin\_K5}: best validation NLL 3.0898, lineage
$[0,2,3,4,9]$, $\alpha\in[0.10009,0.13921]$, median/max alpha-gradient
4.047/20.553, clip fraction 0.8977, runtime 1898.5 s. Финальная $\Omega$
имеет eigenvalue range $[0.0257,17.8022]$; train soft marginal L1 до uniform
равен 0.0969.

\subsection{Трёхсидовая сводка all-12}

\begin{landscape}
\scriptsize
\begin{longtable}{llrrrrr}
\toprule
Dataset & Arch & W test mean$\pm$SD & no-W test mean$\pm$SD & $\Delta$ & W ARI mean & LaL \\
\midrule
\endfirsthead
\toprule
Dataset & Arch & W test mean$\pm$SD & no-W test mean$\pm$SD & $\Delta$ & W ARI mean & LaL \\
\midrule
\endhead
K2\_C5 & COTIC & $1.0000\pm0.0000$ & $1.0000\pm0.0000$ & 0.0000 & 1.0000 & 0.970 \\
K3\_C5 & COTIC & $1.0000\pm0.0000$ & $1.0000\pm0.0000$ & 0.0000 & 1.0000 & 0.850 \\
K4\_C5 & COTIC & $1.0000\pm0.0000$ & $1.0000\pm0.0000$ & 0.0000 & 1.0000 & 0.900 \\
K5\_C5 & COTIC & $1.0000\pm0.0000$ & $0.9817\pm0.0144$ & 0.0183 & 1.0000 & 0.840 \\
sin\_K2\_C5 & COTIC & $1.0000\pm0.0000$ & $1.0000\pm0.0000$ & 0.0000 & 1.0000 & 0.990 \\
sin\_K3\_C5 & COTIC & $0.9917\pm0.0000$ & $0.9889\pm0.0048$ & 0.0028 & 0.9759 & 0.990 \\
sin\_K4\_C5 & COTIC & $0.9458\pm0.0219$ & $0.9563\pm0.0108$ & -0.0104 & 0.8555 & 0.920 \\
sin\_K5\_C5 & COTIC & $0.9883\pm0.0076$ & $0.9867\pm0.0029$ & 0.0017 & 0.9691 & 0.920 \\
trunc\_K2\_C5 & COTIC & $1.0000\pm0.0000$ & $1.0000\pm0.0000$ & 0.0000 & 1.0000 & 1.000 \\
trunc\_K3\_C5 & COTIC & $0.9778\pm0.0127$ & $0.9778\pm0.0048$ & 0.0000 & 0.9352 & 0.960 \\
trunc\_K4\_C5 & COTIC & $0.9896\pm0.0036$ & $0.9813\pm0.0063$ & 0.0083 & 0.9710 & 0.990 \\
trunc\_K5\_C5 & COTIC & $0.9717\pm0.0029$ & $0.9733\pm0.0058$ & -0.0017 & 0.9311 & 0.940 \\
\midrule
K2\_C5 & THP & $1.0000\pm0.0000$ & $1.0000\pm0.0000$ & 0.0000 & 1.0000 & 0.970 \\
K3\_C5 & THP & $1.0000\pm0.0000$ & $1.0000\pm0.0000$ & 0.0000 & 1.0000 & 0.850 \\
K4\_C5 & THP & $0.9958\pm0.0036$ & $0.9958\pm0.0072$ & 0.0000 & 0.9897 & 0.900 \\
K5\_C5 & THP & $0.9867\pm0.0058$ & $0.9933\pm0.0058$ & -0.0067 & 0.9682 & 0.840 \\
sin\_K2\_C5 & THP & $1.0000\pm0.0000$ & $0.9958\pm0.0072$ & 0.0042 & 1.0000 & 0.990 \\
sin\_K3\_C5 & THP & $0.9889\pm0.0048$ & $0.9889\pm0.0048$ & 0.0000 & 0.9681 & 0.990 \\
sin\_K4\_C5 & THP & $0.8958\pm0.0295$ & $0.8917\pm0.0260$ & 0.0042 & 0.7519 & 0.920 \\
sin\_K5\_C5 & THP & $0.8667\pm0.0900$ & $0.9250\pm0.0260$ & -0.0583 & 0.7423 & 0.920 \\
trunc\_K2\_C5 & THP & $1.0000\pm0.0000$ & $1.0000\pm0.0000$ & 0.0000 & 1.0000 & 1.000 \\
trunc\_K3\_C5 & THP & $0.9722\pm0.0048$ & $0.9694\pm0.0048$ & 0.0028 & 0.9192 & 0.960 \\
trunc\_K4\_C5 & THP & $0.9917\pm0.0036$ & $0.9896\pm0.0036$ & 0.0021 & 0.9775 & 0.990 \\
trunc\_K5\_C5 & THP & $0.9717\pm0.0029$ & $0.9717\pm0.0104$ & 0.0000 & 0.9317 & 0.940 \\
\bottomrule
\end{longtable}
\end{landscape}

\begin{center}
\begin{tabular}{lrrrrrr}
\toprule
Arch & W purity & no-W purity & $\Delta$ & W wins/ties/losses & $\Delta$ NLL & Runtime W/no-W \\
\midrule
THP & 0.97245 & 0.97677 & -0.00432 & 9/18/9 & +0.00212 & 7.76/3.49 h \\
COTIC & 0.98874 & 0.98715 & +0.00159 & 11/20/5 & -0.01548 & 8.83/3.23 h \\
\bottomrule
\end{tabular}
\end{center}

В среднем по datasets THP-Wishart выше опубликованного LaL mean на 8 из 12,
равен на одном и ниже на трёх; COTIC-Wishart выше на 10, равен на одном и ниже
на одном. Это не paired reproduction LaL, а масштабная внешняя reference line.
Совокупный recorded runtime 144 runs составляет 23.31 GPU-hours; peak VRAM в
all-12 не мониторился отдельным sampler.

\subsection{Разбор \code{sin\_K5\_C5} по seeds}

\begin{center}
\small
\begin{tabular}{llrrrrr}
\toprule
Arch & Variant & Seed & Val & Test & ARI & Test NLL \\
\midrule
THP & no-W & 0 & 0.920 & 0.910 & 0.7828 & 3.4083 \\
THP & W & 0 & 0.865 & 0.870 & 0.7166 & 3.3949 \\
THP & no-W & 1 & 0.935 & 0.955 & 0.8879 & 3.2421 \\
THP & W & 1 & 0.940 & 0.955 & 0.8877 & 3.2715 \\
THP & no-W & 2 & 0.915 & 0.910 & 0.7859 & 3.2756 \\
THP & W & 2 & 0.830 & 0.775 & 0.6227 & 3.2985 \\
\midrule
COTIC & no-W & 0 & 0.975 & 0.990 & 0.9764 & 3.0296 \\
COTIC & W & 0 & 0.990 & 0.980 & 0.9474 & 3.0326 \\
COTIC & no-W & 1 & 0.980 & 0.985 & 0.9627 & 3.0484 \\
COTIC & W & 1 & 0.990 & 0.990 & 0.9734 & 3.0047 \\
COTIC & no-W & 2 & 0.975 & 0.985 & 0.9627 & 3.0233 \\
COTIC & W & 2 & 0.985 & 0.995 & 0.9866 & 3.0042 \\
\bottomrule
\end{tabular}
\end{center}

\verdict{All-12 показывает не универсальную пользу Wishart, а interaction с
backbone. COTIC устойчиво решает \code{sin\_K5}; THP имеет seed-dependent
collapse, особенно W seed 2. Небольшой NLL gain THP W seed 0 не компенсирует
ухудшение clustering.}

\section{Ранние и промежуточные эксперименты}

\subsection{Published snapshot и 300-epoch retraining}

На 1200 trajectories (split 90/45/1065) были проверены NHP, THP, COTIC и две
component parameterizations. Таблица показывает новый retraining seed 0.

\begin{center}
\small
\begin{tabular}{lrrrr}
\toprule
Модель & no-W NLL & W NLL & W-no-W & W purity \\
\midrule
NHP/output split & 3.5515 & 3.4580 & -0.0935 & 0.936 \\
NHP/lal fixed & 3.5776 & 3.4719 & -0.1057 & 0.923 \\
THP/output split & 3.6236 & 3.5128 & -0.1108 & 0.661 \\
THP/lal fixed & 3.5368 & 3.4969 & -0.0399 & 0.925 \\
COTIC/output split & 3.9159 & 3.6408 & -0.2751 & 0.934 \\
COTIC/lal fixed & 3.8771 & 3.6732 & -0.2039 & 0.709 \\
\bottomrule
\end{tabular}
\end{center}

Wishart улучшил suffix NLL во всех шести comparisons, но THP/output-split
purity упал с опубликованных 0.951 до 0.661, COTIC/lal-fixed --- с 0.892 до
0.709. Внутренняя MC-SD была мала: нестабилен optimizer basin, не evaluation.

\subsection{Signed DGP: moderate против strong}

Трёхсидовый NHP-эксперимент использовал exact signed K3C5 DGP. В таблице
разность signed$-\alpha=0$ по NLL; отрицательное лучше.

\begin{center}
\begin{tabular}{lrrrr}
\toprule
DGP & Signed NLL & Signed$-\alpha=0$ & Purity & ARI \\
\midrule
Zero & $3.55375\pm0.06726$ & $-0.00591\pm0.01049$ & 0.8514 & 0.7155 \\
Moderate & $3.48672\pm0.06376$ & $-0.03394\pm0.01543$ & 0.9531 & 0.8673 \\
Strong & $3.56485\pm0.10071$ & $+0.02086\pm0.02296$ & 0.9336 & 0.8177 \\
\bottomrule
\end{tabular}
\end{center}

Moderate выиграл на всех seeds, strong проиграл на всех. Увеличение model $\nu$
до 30/60/120 не исправило strong DGP. Уменьшение $\eta_\alpha$ монотонно
удерживало $\alpha$ около 0.1 и улучшало NLL, что больше похоже на ограничение
вредной interaction, чем на успешную идентификацию.

\subsection{THP fixed-split, MC и InfoMax}

\begin{center}
\small
\begin{tabular}{lrrrrr}
\toprule
Variant & Val purity & Test purity & ARI & 1-to-1 accuracy & Все 5 мод \\
\midrule
THP no-W & $0.808\pm0.010$ & $0.710\pm0.014$ & 0.507 & $0.649\pm0.016$ & 0/5 \\
WiRE-THP & $0.819\pm0.004$ & $0.729\pm0.007$ & 0.547 & $0.656\pm0.015$ & 0/5 \\
WiRE + InfoMax 5 & $0.810\pm0.010$ & $0.740\pm0.017$ & 0.555 & $0.705\pm0.049$ & 2/5 \\
WiRE + InfoMax 10 & $0.811\pm0.007$ & $0.749\pm0.019$ & 0.543 & $0.723\pm0.052$ & 3/5 \\
\bottomrule
\end{tabular}
\end{center}

Обе plain-модели использовали пять predicted components, но покрывали только
четыре true modes. Train MC samples 2/4/8 дали val purity 0.820/0.820/0.820;
validation MC 8--128 и test MC 16--128 также не меняли purity. Значит,
four-mode collapse не был следствием недостаточного MC.

\subsection{Two-phase COTIC: сильный, но односидовый результат}

На \code{K5\_C5}: K=1 pretrain 300, expansion $1\to10$, frozen encoder,
DPP 0.1, pruning 300--700, затем DPP $\to0$, backbone LR $0\to10^{-5}$ и
$\rho:5\to50$. Получено val/test purity 0.995/1.000, ARI/NMI 1.000. Direct
joint no-W/W controls дали test 0.665/0.825.

\caveat{Это один dataset и один seed. Поздний all-12 использовал иной протокол:
pretrain 10, encoder unfrozen, DPP 0, $\rho=5$. Поэтому результат доказывает
силу двухфазной optimization geometry, а не чистый эффект Wishart.}

\section{Поздние THP sweeps на \code{sin\_K5\_C5}}

\subsection{Frozen backbone LR \texorpdfstring{$\times$}{x} warmup}

25 seed-0 jobs: backbone LR $\{10^{-5},3\cdot10^{-5},10^{-4},3\cdot10^{-4},10^{-3}\}$
$\times$ warmup $\{300,100,50,25,10\}$. Encoder frozen до target $K$, physical
batch 8, accumulation 8, effective batch 64. Validation top points:

\begin{center}
\begin{tabular}{rrrrrr}
\toprule
Backbone LR & Warmup & Val & Val NLL & Test & ARI \\
\midrule
$3\cdot10^{-4}$ & 50 & 0.945 & 3.2431 & 0.910 & 0.7859 \\
$3\cdot10^{-4}$ & 300 & 0.940 & 3.2091 & 0.910 & 0.7857 \\
$10^{-4}$ & 300 & 0.940 & 3.2401 & 0.895 & 0.7495 \\
$10^{-5}$ & 300 & 0.935 & 3.3018 & 0.905 & 0.7782 \\
$3\cdot10^{-4}$ & 10 & 0.930 & 3.2155 & 0.905 & 0.7778 \\
\bottomrule
\end{tabular}
\end{center}

Односидовый best превосходит THP no-W all-12 seed0 validation 0.920, но
descriptive test 0.910 лишь совпадает с no-W. Это показывает улучшение basin
selection на validation, не доказанный test gain.

\subsection{\texorpdfstring{$\eta_\alpha\times$}{eta-alpha x} K=1 pretraining length}

25 jobs с independently built K=1 checkpoints. Лучший validation result:
pretrain 300, $\eta_\alpha=10^{-5}$, val/test 0.935/0.890. При той же длине
$\eta_\alpha=10^{-4}$ дал 0.925/0.885; $10^{-3}$ --- 0.930/0.895.
Короткий pretrain 10--100 в основном снижал val до 0.795--0.830.

\verdict{Длина K=1 pretraining оказалась важнее точной настройки $\eta_\alpha$.
В этом grid нет evidence, что быстрый alpha optimizer восстанавливает
кластеризацию; на лучших точках $\alpha$ остаётся около 0.10--0.22.}

\subsection{Big-protocol \texorpdfstring{$\beta_0=1$}{beta0=1} и
\texorpdfstring{$\eta_\alpha$}{eta-alpha}}

Здесь воспроизведён фактический общий neural optimizer all-12: backbone и heads
всегда $10^{-3}$, exact K=1 checkpoint all-12 reused. Единственные изменения ---
$\beta_0=1$ и $\eta_\alpha$.

\begin{center}
\begin{tabular}{rrrrrr}
\toprule
$\eta_\alpha$ & Val & Val NLL & Test & ARI & Best $\alpha$ \\
\midrule
$10^{-5}$ & 0.910 & 3.0623 & 0.885 & 0.7367 & 0.1041 \\
$10^{-4}$ & 0.905 & 3.0660 & 0.925 & 0.8164 & 0.1378 \\
$10^{-3}$ & 0.900 & 3.0822 & 0.900 & 0.7666 & 0.3716 \\
$3\cdot10^{-4}$ & 0.900 & 3.0889 & 0.895 & 0.7543 & 0.2175 \\
$3\cdot10^{-5}$ & 0.805 & 3.0675 & 0.735 & 0.5468 & 0.1123 \\
\bottomrule
\end{tabular}
\end{center}

Validation rule выбирает $10^{-5}$. После этого для matched-DGP вручную
зафиксировано $10^{-4}$. Это допустимый новый design choice, но его нельзя
обосновывать более высоким test 0.925.

\subsection{Overcomplete backbone LR, частичный sweep}

Из 7 запланированных завершены 5. До target $K$ heads имели LR $10^{-3}$,
backbone --- swept LR; после target backbone возвращался к $10^{-3}$ за 200
steps. Лучшие завершённые точки: $10^{-6}$ val/test 0.870/0.835; frozen 0
--- 0.855/0.845; $3\cdot10^{-6}$ --- 0.840/0.805. Прогон остановлен, поэтому
он не задаёт финальную конфигурацию.

\section{Matched Wishart random-effect DGP}

\subsection{Почему первый DGP был невалиден}

Первый full suite содержал 2000 sequences и имел почти идеальную oracle
block-mass routing accuracy 0.998, но обе обучаемые модели collapsed в один
cluster: val/test purity 0.2, ARI/NMI 0.0. Hidden $W$ был структурирован, но
observable event data недостаточно отражали random effect, а K=1 cloning и
слабая early separation не создавали usable clusters. Suite остановлен после
2/9 jobs.

\subsection{Исправленная генерация}

Новый DGP сохраняет синусоидальный COTIC-like kernel, но вводит block-local
activity. Для каждой trajectory её истинный $K\times C$ block получает
нормированную local activity $a^{loc}$; между блоками используется global
activity $a^{glob}$:
\[
p=\gamma_{in}(B\odot M_{in})a^{loc}
+\gamma_{out}(B\odot M_{out})a^{glob}.
\]
Финальная DGP intensity использует тот же convex transform с
$\alpha_{true}=0.35$. Текущая постановка: $K=C=5$, 400 sequences/cluster в
полном плане, horizon 20, $\nu_{data}=25$, baseline totals
$[0.8,6.5,6.5,16.0,2.3]$, $\gamma_{in}=3$, $\gamma_{out}=0.25$.

Full preflight на 2000 sequences до block-local gain имел mean events 123.592,
oracle routing 0.9935, rate relative SD 0.0200, validation ridge-LDA accuracy
0.710 и all-data k-means ARI 0.7483. Gain-3 debug data на 400 sequences имели
mean events 171.423, rate relative SD 0.0351, ridge-LDA 0.725 и k-means ARI
0.7038. То есть effect стал сильнее, но задача одновременно стала интенсивнее.

\subsection{Gain 1/3/5 и paired результат gain 3}

Preflight DGP summary:
\begin{center}
\begin{tabular}{rrrrr}
\toprule
Within gain & Rate rel. SD & Ridge-LDA acc & k-means ARI & Mean events \\
\midrule
1 & 0.01395 & 0.875 & 0.7864 & 134.87 \\
3 & 0.03510 & 0.725 & 0.7038 & 171.42 \\
5 & 0.05128 & 0.675 & 0.6546 & 212.78 \\
\bottomrule
\end{tabular}
\end{center}

Gain 3, short COTIC, 3 seeds:
\begin{center}
\small
\begin{tabular}{lrrrrrr}
\toprule
Seed & Val no-W & Val W & $\Delta$ val & Test no-W/W & ARI no-W/W & NLL no-W/W \\
\midrule
0 & 0.775 & 0.900 & +0.125 & 0.675/0.775 & 0.5196/0.6178 & 3.1214/3.1234 \\
1 & 0.825 & 0.800 & -0.025 & 0.725/0.750 & 0.5842/0.5585 & 3.1665/3.0966 \\
2 & 0.875 & 0.825 & -0.050 & 0.750/0.800 & 0.6040/0.6519 & 3.1119/3.1026 \\
\midrule
Mean & 0.825 & 0.8417 & +0.0167 & 0.7167/0.7750 & 0.5692/0.6094 & 3.1333/3.1076 \\
\bottomrule
\end{tabular}
\end{center}

\verdict{Гипотеза получила только слабое exploratory evidence: средние purity,
ARI и NLL лучше у Wishart, но validation sign положителен лишь на одном из трёх
seeds. Этого недостаточно для запуска дорогого full 9-task suite как
confirmatory benchmark.}

\section{Обучаемая \texorpdfstring{$\nu$}{nu} и Monte Carlo}

\subsection{Непрерывная параметризация \texorpdfstring{$\nu$}{nu}}

Для learnable mode реализовано
\[
\nu=D+(\nu_{max}-D)\sigma(r_\nu),\qquad D<\nu<\nu_{max},
\]
и weak log-scale prior
\[
\mathcal L_\nu=\frac{\gamma_\nu}{2}
\left(\log\frac{\nu}{\nu_0}\right)^2.
\]
Wishart sample строится дифференцируемым Bartlett factor: нижнетреугольные
off-diagonal элементы standard normal, диагональ
$\sqrt{\chi^2_{\nu-j+1}}$, реализованная через reparameterized gamma. Fixed
integer branch оставлен byte-compatible.

\rationale{Bounded sigmoid гарантирует $\nu>D$, а Bartlett decomposition
разрешает gradient по continuous degrees of freedom. Отдельный optimizer group
и слабый prior предотвращают мгновенный уход к границе.}

\subsection{Fixed 25, fixed 37 и learned 37}

\begin{center}
\begin{tabular}{lrrrrr}
\toprule
Mode & Val purity & Val NLL & Test purity & Test ARI & Selected $\nu$ \\
\midrule
Fixed 25 & 0.8583 & 3.10216 & 0.7667 & 0.6013 & 25 \\
Fixed 37 & 0.8500 & 3.10753 & 0.7750 & 0.6094 & 37 \\
Learned from 37 & 0.8583 & 3.10600 & 0.7417 & 0.5684 & 37.505 mean \\
\bottomrule
\end{tabular}
\end{center}

По validation rule fixed 25 выигрывает tie по purity меньшим NLL. Learned $\nu$
по seeds выбрал 37.232, 37.809, 37.475; gradients были малы, параметр остался
рядом со стартом.

\subsection{\texorpdfstring{LR$_\nu$}{LR-nu} и train-MC samples}

Seed 0: $\eta_\nu=10^{-3}$ дал val 0.850 и $\nu=37.23$;
$3\cdot10^{-3}$ --- 0.825 и 37.84; $10^{-2}$ --- 0.825 и 40.78. Больший LR
сдвигает $\nu$ вверх, но purity ухудшается.

\begin{center}
\begin{tabular}{rrrrrrr}
\toprule
$S_{train}$ & Val purity & Val NLL & Selected $\nu$ & $|g_\nu|_{max}$ & Runtime & Peak VRAM \\
\midrule
1 & 0.825 & 3.12222 & 37.274 & 0.01343 & 30.6 s & 1803 MiB \\
2 & 0.825 & 3.12380 & 37.531 & 0.00893 & 32.8 s & 2029 MiB \\
4 & 0.800 & 3.10675 & 36.985 & 0.00435 & 35.8 s & 2921 MiB \\
\bottomrule
\end{tabular}
\end{center}

\verdict{Больше MC samples действительно уменьшает gradient noise и при
$S=4$ улучшает validation NLL, но primary clustering metric падает. Проблема не
объясняется только MC variance.}

\section{Обоснование составных частей: итоговая матрица}

\begin{longtable}{p{0.20\textwidth}p{0.33\textwidth}p{0.37\textwidth}}
\toprule
Составная часть & Теоретическая роль & Экспериментальный статус \\
\midrule
\endfirsthead
\toprule
Составная часть & Теоретическая роль & Экспериментальный статус \\
\midrule
\endhead
K=1 pretrain & Отделить density representation от cluster allocation. & Сильный practical effect; слишком короткий pretrain ухудшает THP. \\
Clone + noise & Сохранить scale и разрушить symmetry. & Необходим; неправильная cloning scale приводила к collapse matched-DGP. \\
Overcomplete $K_0$ & Дать elimination запас кандидатов. & Работает на простых DAN; не гарантирует все true modes на sin-K5 THP. \\
Shared $W_i$ & Связать routing и trajectory-level intensity effect. & NLL benefit встречается часто; clustering benefit backbone- и seed-dependent. \\
Correlation $R$ и $B=R-I$ & Отделить связи от scale и self-effect. & Численно стабильно; само по себе не обеспечивает идентификацию. \\
Projected $\alpha$ & Ограничить interaction и сохранить exact no-effect point. & Большой $\eta_\alpha$ часто вреден; optimum обычно близок к 0.1. \\
Exploration $\beta_t$ & Отложить ранний basin lock-in. & Средний эффект мал; beta$_0=1$ не дал устойчивой победы. \\
Global OT dual & Препятствовать component collapse глобально. & Критически важен; weak/incorrect setting collapsed matched-DGP. \\
Hyperprior $\Omega$ & SPD barrier и shrinkage. & Предотвращает вырождение; learned matrices всё же могут быть анизотропны. \\
Functional DPP & Развести клонированные heads до pruning. & Сильный в one-seed two-phase COTIC; factorized multiseed evidence отсутствует. \\
Physical elimination & Удалять least harmful head без labels. & Даёт auditable lineage; дорого по runtime, чувствительно к representation drift. \\
Encoder freeze & Сделать соседние pruning decisions сопоставимыми. & Улучшил отдельные THP/COTIC runs; all-12 его не использовал. \\
Gradient accumulation & Стабилизировать global marginal при ограниченной VRAM. & Позволяет effective batch 128; memory-safe variants сохраняют семантику. \\
MC 2/8/16/8 & Баланс variance и стоимости. & Увеличение train/eval MC purity не улучшило; retained default оправдан. \\
Learnable $\nu$ & Адаптировать variance random effect. & Реализуемо и стабильно, но в debug не улучшило validation-selected clustering. \\
Validation-only selection & Исключить test leakage. & Обязательный инвариант; ручная фиксация после test не считается подтверждением. \\
\bottomrule
\end{longtable}

\section{Вычислительная среда и memory safety}

Исходная реконструкция: Windows, Python 3.12.13, PyTorch 2.12.0+cu126,
NumPy 2.3.5, pandas 2.2.3, EasyTPP 0.2.1, RTX 4060 Laptop 8188 MiB.
All-12 использовал physical batch 16/64 и effective batch 128 без отдельного
VRAM sampler. Поздние sweeps были закреплены за GPU 0. Они использовали physical
batch 8, pruning batch 2, accumulation 16 и guard при sampled dedicated VRAM
$\ge7000$ MiB.

После 06.08 physical-removal scoring сначала генерирует прежнюю RNG-группу
Wishart samples, затем строит THP traces адаптивными microbatches с budget не
более 512 padded events; candidate scores и common-random-number semantics
сохраняются, CUDA cache очищается между candidate copies.

\rationale{Memory adaptation меняет только физическое разбиение. Effective
batch, Wishart samples и candidate RNG остаются прежними, поэтому сравнение
scores не превращается в другую статистическую процедуру.}

\section{Ошибки протокола, исправления и источники неопределённости}

\begin{enumerate}
  \item \textbf{Metadata LR mismatch.} Config all-12 содержит backbone LR
  $10^{-5}$, history --- $10^{-3}$. Фактический history имеет приоритет.
  \item \textbf{Original seed-0 Wishart rows.} В старом suite использовались
  прежние W settings; в основной таблице они заменены corrected 24/24.
  \item \textbf{Неидеальный seed-0 pairing.} Corrected W и historical no-W
  имеют разные K=1 checkpoints. Seeds 1--2 paired byte-identically.
  \item \textbf{Different no-W dual LR.} Значит, all-12 W-noW --- сравнение
  исторических variants, не чистая абляция одного слоя.
  \item \textbf{Single-seed tuning.} Большинство Aug6--7 grids оптимизировалось
  на seed 0; их ranking нельзя считать multiseed generalization.
  \item \textbf{Multiple exploration.} После многих grids validation metrics
  сами становятся объектом adaptive search; нужен fresh confirmatory seed set.
  \item \textbf{LaL scope.} Published means --- внешняя reference; локальный
  BOS-control неполон и не заменяет LaL random walk.
\end{enumerate}

\section{Рекомендуемый следующий подтверждающий эксперимент}

Новый дорогой sweep сейчас не оправдан. Если работу продолжать, минимальный
научно чистый протокол должен быть следующим:

\begin{enumerate}
  \item Зафиксировать один matched gain-3 DGP generator до обучения и сохранить
  data hashes; не менять DGP после просмотра metrics.
  \item Сравнить только no-W и W с exact shared K=1 checkpoint, одинаковыми
  neural/dual LRs, batch/RNG schedule и без test-guided tuning.
  \item Зафиксировать W: $\eta_\alpha=10^{-4}$ как заранее объявленный design,
  $\nu=25$ fixed, $S_{train/val/test}=2/8/16$, $\beta_0=1$.
  \item Использовать минимум 10 fresh model seeds; primary endpoint --- paired
  validation purity W-noW, secondary --- ARI/NMI/NLL. Test открыть один раз
  после заморозки анализа.
  \item Отдельно, а не в том же comparison, проверить two-phase freeze+DPP,
  чтобы не смешивать effect Wishart с effect optimization geometry.
  \item Полный LaL comparison называть таковым только после реализации его
  random-walk/split-delete protocol; до этого использовать термин
  ``LaL-style fixed-K BOS control''.
\end{enumerate}

\section{Артефакты и воспроизводимость}

\subsection{Главные источники all-12}
\sourcepath{artifacts/dan_all12_best_elimination_corrected_seed0_20260805/summary.csv}
\sourcepath{artifacts/dan_all12_best_elimination_seed0_20260804/summary.csv}
\sourcepath{artifacts/dan_all12_best_elimination_seeds1_2_20260805/summary.csv}
\sourcepath{scripts/run_two_phase_all12_benchmark.py}
\sourcepath{artifacts/dan_all_arch_newwishart_beta_cotic7_seed0_20260802/PUBLISHED_DAN_PURITY.csv}

Для точного THP \code{sin\_K5} источниками являются
\code{config.json}, \code{history.csv}, \code{progress.json},
\code{pruning\_events.json} и \code{result.json} внутри соответствующего
run directory. При конфликте фактического LR используются history/progress.

\subsection{Отчёты предшествующих фаз}
\sourcepath{REPRODUCTION_REPORT_20260801.md}
\sourcepath{SIGNED_EXPERIMENT_REPORT_20260801.md}
\sourcepath{THP_WISHART_OPTUNA_EXPLORATION_REPORT_20260802.md}
\sourcepath{THP_WISHART_SPLIT42_EPOCHWISE_REPORT_20260802.md}
\sourcepath{COTIC_TWO_PHASE_WIRE_REPORT_20260803.md}

\subsection{Поздние sweeps и matched DGP}
\sourcepath{artifacts/dan_sin_k5_thp_frozen_backbone_lr_warmup_seed0_20260806/sweep_summary.csv}
\sourcepath{artifacts/dan_sin_k5_thp_frozen_lr_alpha_pretrain_seed0_20260806/sweep_summary.csv}
\sourcepath{artifacts/dan_sin_k5_thp_big_protocol_beta1_lr_alpha_seed0_20260807/sweep_summary.csv}
\sourcepath{configs/experiments/wishart_re_sin_k5_cotic_comparison.json}
\sourcepath{artifacts/_debug_wishart_re_blocklocal_gain3_preflight_20260807/dgp_audit.json}
\sourcepath{artifacts/_debug_wishart_re_blocklocal_gain3_preflight_20260807/summary.csv}
\sourcepath{artifacts/_debug_wishart_re_gain3_nu_fixed25_20260807/summary.csv}
\sourcepath{artifacts/_debug_wishart_re_gain3_nu_learned37_20260807/summary.csv}
\sourcepath{artifacts/_debug_wishart_re_gain3_nu_learned37_mc4_val8_20260807/summary.csv}

\subsection{Реализация алгоритма}
\sourcepath{src/lal_wishart/models/wishart_math.py}
\sourcepath{src/lal_wishart/models/signed_wishart.py}
\sourcepath{src/lal_wishart/models/latent_wishart_attention_nhp.py}
\sourcepath{src/lal_wishart/train/fit_latent_wishart_attention_nhp.py}
\sourcepath{src/lal_wishart/reproduction/wishart_re_sin_k5c5.py}
\sourcepath{scripts/run_wishart_re_sin_k5_cotic_comparison.py}

\section{Финальный вывод}

\verdict{Главный положительный результат проекта --- не универсальный clustering
gain, а воспроизводимая и технически целостная процедура trajectory-level
Wishart marginalization, physical elimination и validation-only selection.
Predictive NLL нередко улучшается. Однако clustering advantage зависит от
backbone, initialization и DGP; на THP \code{sin\_K5} Wishart может ухудшать
результат, а matched-DGP gain остаётся нестабильным по seeds. Поэтому текущий
научно корректный вывод --- отрицательный/неопределённый, а не ``Wishart лучше''.
Следующий шаг должен быть один заранее зарегистрированный paired multiseed
эксперимент, а не очередной seed-0 sweep.}

\end{document}
